% ============================================================================
% 04_capitulo4_metodologia.tex
% Reescrita e complementos do Capitulo 4 -- METODOLOGIA.
%
% Como usar:
%   (A) Inserir a nova Secao 4.1.3 (Enriquecimento via APIs externas)
%       como ultima subsecao da Secao 4.1 (Coleta de Dados).
%   (B) Adicionar o paragrafo do bloco 4.B ao FINAL da Secao 4.2
%       (Tecnologias, Bibliotecas e Ferramentas).
%   (C) Adicionar as Secoes 4.3.7 e 4.3.8 ao FINAL da Secao 4.3
%       (Pre-processamento dos dados).
%   (D) SUBSTITUIR INTEIRAMENTE a Secao 4.4 (Treinamento dos Modelos
%       de Previsao) pelo bloco 4.D.
%   (E) SUBSTITUIR/EXPANDIR a Secao 4.5 (Avaliacao) pelo bloco 4.E,
%       incorporando interpretabilidade, threshold tuning e validacao
%       temporal.
%   (F) SUBSTITUIR a Secao 4.6 (Geracao do Mapa de Risco) pelo bloco
%       4.F, que descreve a aplicacao web interativa explicavel.
%   (G) Adicionar a nova Secao 4.7 (Pseudo-labeling de Umidade) ao
%       FINAL do Capitulo 4.
% ============================================================================

% ============================================================================
% (A) NOVA SUBSECAO 4.1.3 -- Enriquecimento via APIs externas
% ============================================================================

\subsection{Enriquecimento via APIs externas (NASA POWER e NASA FIRMS)}
\label{subsec:metod_apis_externas}

Além do Programa Queimadas, duas APIs externas são consultadas para enriquecer o conjunto de dados:

\paragraph{NASA POWER (\emph{Prediction of Worldwide Energy Resources})~\cite{nasapower}.} Serve reanálise climática global em resolução de \(0{,}5^\circ \times 0{,}625^\circ\) com latência diária. Foi utilizada para extrair a variável RH2M (umidade relativa a 2~m), reconhecida na literatura como preditor relevante de fogo. A integração foi feita via \texttt{scripts/enriquecer\_dados\_umidade.py}, com \emph{cache} em disco, retomada após interrupção e respeito ao limite de taxa (\(\approx 32\)~requisições/min com duas chaves de API). Pela alta latência de enriquecimento total (\(\approx 16\) dias para \num{755711} requisições únicas), o projeto adotou estratégia complementar: \num{170465} registros foram efetivamente enriquecidos (\(\approx \SI{18,3}{\percent}\) da base) e os \num{763489} restantes foram imputados via \emph{pseudo-labeling} com LightGBM \emph{regressor} (\(R^2 = 0{,}933\) em \emph{holdout} --- Seção~\ref{sec:metodologia_pseudo_label}).

\paragraph{NASA FIRMS (\emph{Fire Information for Resource Management System})~\cite{nasafirms}.} Serve detecções ativas de fogo em quase-tempo-real, com FRP (\emph{Fire Radiative Power}) por píxel para os sensores VIIRS e MODIS. É consultada em \textbf{tempo de inferência} pelo aplicativo web (Seção~\ref{sec:metod_app_web}) para incrementar a \emph{feature} \texttt{FRP} com observações dos últimos 7 dias na vizinhança do ponto consultado.

% ============================================================================
% (B) COMPLEMENTO PARA A SECAO 4.2 -- Tecnologias, Bibliotecas e Ferramentas
% ============================================================================

\bigskip
\noindent
Além das bibliotecas já listadas (\emph{Pandas}~\cite{mckinney2010data}, \emph{NumPy}~\cite{harris2020numpy}, \emph{Scikit-Learn}~\cite{pedregosa2011scikit}, \emph{Matplotlib}, \emph{Seaborn}, \emph{Joblib}, \emph{Geopy}), o \emph{pipeline} atual incorpora:

\begin{itemize}
  \item \textbf{XGBoost} (\texttt{xgboost})~\cite{chen2016xgboost}, \textbf{LightGBM} (\texttt{lightgbm})~\cite{ke2017lightgbm} e \textbf{CatBoost} (\texttt{catboost==1.2.10})~\cite{prokhorenkova2018catboost} --- modelos de \emph{gradient boosting} usados como classificadores autônomos e como \emph{base learners} dos \emph{ensembles}.
  \item \textbf{Imbalanced-learn} (\texttt{imblearn})~\cite{lemaitre2017imbalanced} --- implementação do SMOTE para a variante \texttt{random\_forest\_smote}.
  \item \textbf{Optuna} (\texttt{optuna})~\cite{akiba2019optuna} --- otimização bayesiana de hiperparâmetros via TPE.
  \item \textbf{SHAP} (\texttt{shap}) --- interpretabilidade local e global via \texttt{TreeExplainer}~\cite{lundberg2017unified,lundberg2020local}.
  \item \textbf{Flask} (\texttt{flask}) --- servidor web para a aplicação interativa.
  \item \textbf{Folium} (\texttt{folium})~\cite{folium} --- visualização cartográfica baseada em \emph{Leaflet}~\cite{leafletjs} integrada à renderização do Flask.
  \item \textbf{Geopandas} (\texttt{geopandas})~\cite{geopandas} e \textbf{Shapely} (\texttt{shapely}) --- manipulação de geometrias e do \emph{shapefile} da Amazônia Legal~\cite{ibgeamazonia2024}.
  \item \textbf{Requests} (\texttt{requests}) --- chamadas HTTP às APIs NASA POWER e NASA FIRMS.
\end{itemize}

% ============================================================================
% (C) NOVAS SUBSECOES 4.3.7 e 4.3.8 -- ENGENHARIA DE FEATURES e CALIBRACAO
% ============================================================================

\subsection{Engenharia de \emph{features} físico-climáticas (Tier 1)}
\label{subsec:metod_features_tier1}

Em complemento às nove \emph{features} originais do Programa Queimadas, foram derivadas \textbf{24~\emph{features} físico-climáticas avançadas} (Tier~1), agrupadas em sete camadas (implementação em \texttt{scripts/features\_avancadas.py}). A Tabela~\ref{tab:metod_tier1} resume as camadas com a justificativa bibliográfica de cada uma.

\begin{table}[h]
\centering
\caption{Camadas Tier 1 de \emph{features} físico-climáticas adicionadas ao \emph{dataset}.}
\label{tab:metod_tier1}
\small
\begin{tabular}{clp{6cm}l}
\toprule
\textbf{\#} & \textbf{Camada} & \textbf{\emph{Features}} & \textbf{Justificativa} \\
\midrule
1 & Médias móveis estendidas & \texttt{Precipitacao\_ma14/30/90}, \texttt{DiaSemChuva\_ma14/30/90} & \emph{Fuel moisture lag}~\cite{seager2015climatology} \\
2 & Precipitação acumulada & \texttt{Precipitacao\_acum\_30/90/180/365d} & \cite{seager2015climatology} \\
3 & SPI & \texttt{SPI\_1m}, \texttt{SPI\_3m}, \texttt{SPI\_6m} & \cite{mckee1993spi,npjhazards2025fire} \\
4 & Anomalia de precipitação & \texttt{Anomalia\_Precipitacao}, \texttt{Anomalia\_Precipitacao\_rel} & \cite{forests2024cerradoamazon} \\
5 & KBDI, VPD, De Martonne & \texttt{Temp\_Climatologica}, \texttt{KBDI\_proxy}, \texttt{VPD\_proxy}, \texttt{Aridez\_DeMartonne} & \cite{keetch1968drought,seager2015climatology,demartonne1926aridity,forests2024cerradoamazon} \\
6 & Histórico estendido de incêndios & \texttt{Incendios\_Ultimos\_90/180/365\_Dias}, \texttt{Media\_FRP\_Celula\_30d} & \cite{cheerala2025probabilistic} \\
7 & Estação seca acumulada & \texttt{Dias\_Secos\_90d} & \emph{Canadian Drought Code} (FWI) \\
\bottomrule
\end{tabular}
\end{table}

A granularidade espacial das camadas 1, 2, 6, 7 e do SPI usa \textbf{células de \(0{,}25^\circ\)} (\(\approx \SI{27}{\km}\)), o que evita colapso das janelas temporais que ocorreria com agrupamento por coordenadas literais (cada foco é um ponto único). A \emph{climatologia} de temperatura provém de uma tabela estática de médias mensais por estado, aproximada a partir das Normais Climatológicas do INMET 1981--2010~\cite{inmetnormais}. O CSV final \texttt{base\_de\_dados\_enriquecido.csv} contém \num{924306} linhas e 46~colunas, com tempo total de geração de \(\approx \SI{40}{\second}\) sobre os \num{933954} registros brutos.

\subsection{Calibração de probabilidades}
\label{subsec:metod_calibracao}

A predição padrão dos modelos retorna escores não-calibrados. Para que a saída do modelo seja interpretável como ``\(P(\text{classe} \mid \text{features}) = 0{,}73\)'', aplicou-se \texttt{CalibratedClassifierCV(method='isotonic', cv=3)} aos modelos do \emph{ensemble} final. A calibração isotônica~\cite{zadrozny2002transforming} ajusta uma função monotônica não-paramétrica entre escores brutos e probabilidades empíricas, sem assumir forma sigmoide --- escolha defensável dada a natureza não linear dos GBMs envolvidos.

% ============================================================================
% (D) SECAO 4.4 -- REESCRITA INTEGRAL (TREINAMENTO DOS MODELOS)
% ============================================================================

\section{Treinamento dos modelos de previsão}
\label{sec:metod_treinamento}

\subsection{Divisão dos dados e construção do \emph{pipeline}}

O conjunto de dados foi dividido em treino (\SI{80}{\percent}) e teste (\SI{20}{\percent}) utilizando \texttt{train\_test\_split(stratify=y, random\_state=42)}, preservando a distribuição das três classes em ambos os subconjuntos. Os metadados do \emph{split} (\texttt{random\_state}, \texttt{test\_size}, estratificação) ficam persistidos em \texttt{modelos/split\_metadata.json} para garantir reprodutibilidade. O conjunto de teste contém \num{184862} observações.

O \emph{pipeline} de processamento foi implementado via \texttt{sklearn.pipeline.Pipeline}~\cite{pedregosa2011scikit}, composto por:

\begin{enumerate}
  \item \texttt{ColumnTransformer} --- aplica \texttt{SimpleImputer(strategy='median')} + \texttt{StandardScaler} às \emph{features} numéricas e \texttt{SimpleImputer(strategy='most\_frequent')} + \texttt{OneHotEncoder} às \emph{features} categóricas.
  \item \texttt{CalibratedClassifierCV(method='isotonic', cv=3)} --- envolve o classificador final para garantir calibração das probabilidades.
  \item \textbf{Estimador} --- variável conforme o modelo avaliado.
\end{enumerate}

\subsection{Modelos avaliados}

Oito abordagens foram comparadas no mesmo conjunto de treino/teste:

\begin{enumerate}
  \item \textbf{Regressão Logística balanceada} --- \texttt{class\_weight='balanced'}, hiperparâmetros otimizados via Optuna (\texttt{C=12.43, max\_iter=4768}).
  \item \textbf{SGDClassifier} --- perda logística, \emph{online learning}; \emph{baseline} simples original do projeto.
  \item \textbf{Random Forest com SMOTE}~\cite{chawla2002smote} --- variante para tratamento de desbalanceamento amostral.
  \item \textbf{Random Forest balanceado otimizado}~\cite{breiman2001random} --- \texttt{class\_weight='balanced'}, Optuna (15 \emph{trials}, 3-\emph{fold} CV, \(\text{F}_1\)-macro): \texttt{n\_estimators=346, max\_depth=25, min\_samples\_split=3, max\_features='sqrt'}.
  \item \textbf{XGBoost otimizado}~\cite{chen2016xgboost} --- Optuna (15 \emph{trials}, 3-\emph{fold} CV, \(\text{F}_1\)-macro, \emph{subsample} de 250\,k): \texttt{n\_estimators=346, learning\_rate=0.073, max\_depth=14, subsample=0.74, colsample\_bytree=0.93, gamma=0.03}.
  \item \textbf{LightGBM otimizado}~\cite{ke2017lightgbm} --- Optuna (15 \emph{trials}, \(\text{F}_1\)-macro): \texttt{n\_estimators=683, learning\_rate=0.066, num\_leaves=204, min\_child\_samples=24, subsample=0.78, colsample\_bytree=0.77}.
  \item \textbf{CatBoost}~\cite{prokhorenkova2018catboost} --- \texttt{iterations=1000, learning\_rate=0.05, depth=8, auto\_class\_weights='Balanced'}.
  \item \textbf{Ensemble Voting Soft} --- média das probabilidades calibradas de RF, LR, XGBoost, LightGBM e CatBoost.
  \item \textbf{Ensemble Stacking}~\cite{wolpert1992stacked} --- RF, LR, XGBoost, LightGBM e CatBoost como \emph{base learners}; \texttt{GradientBoostingClassifier}~\cite{friedman2001greedy} como meta-classificador (\texttt{n\_estimators=200, max\_depth=3, learning\_rate=0.1}). Esta é a configuração final adotada (\texttt{ensemble\_stacking\_gbm}).
\end{enumerate}

Todos os modelos foram salvos via \texttt{joblib.dump()} em arquivos \texttt{.pkl} no diretório \texttt{modelos/}. Os hiperparâmetros otimizados via Optuna ficam persistidos em \texttt{modelos/relatorios/optuna\_*.json}.

\subsection{Tempo de treinamento}

A Tabela~\ref{tab:metod_tempos} apresenta o tempo de treinamento de cada abordagem em ambiente CPU com 8~núcleos e \SI{16}{\giga\byte} de RAM.

\begin{table}[h]
\centering
\caption{Tempo de treinamento por modelo.}
\label{tab:metod_tempos}
\small
\begin{tabular}{lr}
\toprule
\textbf{Modelo} & \textbf{Tempo (segundos)} \\
\midrule
Regressão Logística balanceada & 25 \\
SGDClassifier                  & 8 \\
Random Forest + SMOTE          & 480 \\
Random Forest balanceado (Optuna) & 1\,100 \\
XGBoost (Optuna, 250k \emph{subsample}) & 540 \\
LightGBM (Optuna, 250k \emph{subsample}) & 290 \\
CatBoost                       & 410 \\
Ensemble Voting Soft           & 2\,360 (soma dos \emph{base learners} + agregação) \\
\textbf{Ensemble Stacking GBM} & \textbf{5\,100 (\(\approx\) 85 min)} \\
\bottomrule
\end{tabular}
\end{table}

% ============================================================================
% (E) SECAO 4.5 -- EXPANSAO DA AVALIACAO
% ============================================================================

\section{Avaliação dos modelos de previsão}
\label{sec:metod_avaliacao}

A avaliação foi realizada sobre o conjunto de teste de \num{184862} observações. Quatro componentes complementares foram considerados: métricas tradicionais (Seção~\ref{subsec:metod_metricas}), interpretabilidade (Seção~\ref{subsec:metod_interpretabilidade}), ajuste multi-classe de limiares (Seção~\ref{subsec:metod_thresholds}) e validação temporal estrita (Seção~\ref{subsec:metod_validacao_temporal}).

\subsection{Métricas utilizadas}
\label{subsec:metod_metricas}

\begin{itemize}
  \item \textbf{Acurácia global} --- fração de predições corretas (\texttt{accuracy\_score}).
  \item \textbf{\(\text{F}_1\) por classe e \(\text{F}_1\)-macro} --- \texttt{f1\_score(average='macro')}. A \(\text{F}_1\)-macro é a média não ponderada entre as três classes; é a métrica primária por ser robusta a desbalanceamento.
  \item \textbf{Precisão e revocação por classe} --- via \texttt{classification\_report}.
  \item \textbf{Matriz de confusão normalizada} --- \texttt{confusion\_matrix(normalize='true')}.
  \item \textbf{Confiança média e distribuição} --- média e histograma de \(\max\) de \texttt{predict\_proba(X\_test)}.
  \item \textbf{Curva de calibração} --- \texttt{calibration\_curve}, para confirmar que a probabilidade predita corresponde à frequência empírica observada.
\end{itemize}

\subsection{Análise de interpretabilidade}
\label{subsec:metod_interpretabilidade}

Duas técnicas complementares são aplicadas para identificar as \emph{features} mais relevantes:

\paragraph{SHAP por classe via Random Forest leve \emph{proxy}~\cite{lundberg2017unified,lundberg2020local}.} O \texttt{TreeExplainer} aplicado ao \emph{Stacking GBM} final e ao RF Optuna produtivo demanda mais de \num{2}\,\text{GiB} por classe em CPU pós-OHE. Como solução cientificamente defensável~\cite{lundberg2020local}, treina-se um Random Forest \emph{compacto} (100 árvores, \texttt{max\_depth=12}, \texttt{class\_weight='balanced'}, \num{60000} amostras de treino) como \emph{proxy interpretativo}. O script \texttt{scripts/analise\_shap\_rf\_leve.py} aplica \texttt{shap.TreeExplainer} sobre 500 amostras de teste e exporta as importâncias médias absolutas por classe em \texttt{modelos/relatorios/shap\_per\_class.json}.

\paragraph{\emph{Permutation Importance}~\cite{breiman2001random,fisher2019all}.} Aplicada diretamente sobre o \texttt{ensemble\_stacking\_gbm} em produção (5\,000 amostras de teste, três repetições, \emph{seed}=42). Para cada \emph{feature}, calcula-se a queda em \(\text{F}_1\) por classe e em \(\text{F}_1\)-macro quando os valores são permutados aleatoriamente. Diferentemente da importância Gini/MDI, é \emph{model-agnostic} e \textbf{não tendenciosa} para \emph{features} de alta cardinalidade~\cite{strobl2007bias} --- característica relevante dado o OHE de \texttt{Municipio} (542 categorias). A saída fica em \texttt{modelos/relatorios/permutation\_importance\_por\_classe.json}.

\subsection{Ajuste multi-classe de limiares de decisão}
\label{subsec:metod_thresholds}

A predição padrão usa \texttt{argmax} sobre as probabilidades calibradas. Embora maximize \(P(\text{classe} \mid x)\), penaliza a classe \emph{Moderado} --- fronteira ambígua entre \emph{Baixo} e \emph{Muito Alto}~\cite{he2009learning}. O procedimento implementado em \texttt{scripts/ajustar\_threshold.py} varre uma grade de pares (\texttt{thr\_Moderado}, \texttt{thr\_Muito\_Alto}) \(\in \{0{,}30; 0{,}35; \ldots; 0{,}55\}^2\) (36 combinações), com critério primário \(\text{F}_1\)-macro e secundário \(\text{F}_1\)-Moderado. Os limiares finais são persistidos em \texttt{modelos/prediction\_thresholds.json} e aplicados em \emph{runtime} no aplicativo web via parâmetro \texttt{estrategia\_thresholds} do \emph{endpoint} \texttt{/api/predict}.

A lógica de decisão multi-classe em \emph{runtime} é:

\begin{verbatim}
se      P(Moderado)   >= thr_Moderado    -> classe = "Moderado"
senao se P(Muito Alto) >= thr_Muito_Alto -> classe = "Muito Alto"
senao                                    -> classe = "Baixo"
\end{verbatim}

\subsection{Validação temporal (\emph{rolling-origin})}
\label{subsec:metod_validacao_temporal}

Em complemento à validação no \emph{split} aleatório, foi implementado em \texttt{scripts/validacao\_temporal.py} um procedimento de \textbf{\emph{rolling-origin} por ano}~\cite{bergmeir2012cv}. Para cada ano \(Y \in \{2021, 2022, 2023\}\) (três últimos com massa suficiente), o modelo é treinado em \(\text{Ano} < Y\) (\emph{subsample} estratificado de \num{60000} amostras por \emph{fold} para o \emph{Stacking}, limite de memória do OHE denso de \texttt{Municipio}) e avaliado em \(\text{Ano} = Y\). Cobre o uso operacional realista: ``treine com o passado, prediga o ano seguinte''. O resultado é apresentado no Capítulo~5.

% ============================================================================
% (F) SECAO 4.6 -- APLICACAO WEB INTERATIVA EXPLICAVEL
% ============================================================================

\section{Aplicação web interativa explicável}
\label{sec:metod_app_web}

Para tornar o modelo final utilizável e auditável por gestores ambientais, foi desenvolvida uma aplicação web (\texttt{scripts/app\_map\_interativo.py}) baseada em \emph{Flask} + \emph{Folium}~\cite{folium} + \emph{Leaflet}~\cite{leafletjs} que serve como interface explicável sobre o \emph{pipeline}. O usuário clica em qualquer ponto da Amazônia Legal e o sistema responde, em \(\approx 10\)--\SI{20}{\second}, com a predição classificada, acompanhada das 24~\emph{features} Tier~1 calculadas para o ponto, da explicação local das \emph{features} mais influentes e do contexto histórico comparativo.

\subsection{\emph{Pipeline} de predição pontual}

Como o aplicativo prediz \textbf{um único ponto isolado}, sem histórico temporal local, as \emph{features} Tier~1 que dependem de séries (SPI, acumulados longos, histórico estendido de focos) não podem ser calculadas em tempo real. Adota-se a estratégia de \textbf{\emph{lookup} espaço-sazonal} (\texttt{scripts/feature\_lookup.py}):

\begin{enumerate}
  \item Pré-cálculo das medianas das 24~\emph{features} Tier~1 do \emph{dataset} enriquecido, agregadas em três níveis de granularidade decrescente: \((\text{LatBin}~0{,}25^\circ, \text{LonBin}~0{,}25^\circ, \text{Mes})\), depois \((\text{Estado}, \text{Mes})\) e, por fim, \text{Mes} global.

  \item Para o ponto consultado, \emph{lookup} em cascata (do mais fino ao mais grosso) com tolerância de \(\pm 1\) célula.

  \item \emph{Features} físicas independentes de série temporal (\texttt{Temp\_Climatologica}, \texttt{KBDI\_proxy}, \texttt{VPD\_proxy}, \texttt{Aridez\_DeMartonne}, \texttt{Anomalia\_Precipitacao}) são recalculadas em tempo real com o clima NASA POWER atual e a climatologia local da precipitação --- o \emph{lookup} fornece apenas os ``blocos de construção'' históricos imutáveis.
\end{enumerate}

\subsection{Integração de dados em tempo real}

\begin{itemize}
  \item \textbf{Geocodificação reversa} via Nominatim/OpenStreetMap para extrair Estado e Município do clique.
  \item \textbf{Clima atual} via NASA POWER~\cite{nasapower} (precipitação, dias sem chuva, RH2M dos últimos 28~dias).
  \item \textbf{Focos ativos} via NASA FIRMS~\cite{nasafirms} para \texttt{FRP} dos últimos 7~dias no raio de \SI{10}{\km} do ponto.
\end{itemize}

\subsection{Explicabilidade local}

Em \emph{runtime}, \texttt{scripts/explainer.py} aproxima a contribuição de cada \emph{feature} como
\begin{equation}
\mathrm{contrib}(f) = \mathrm{importância}(f) \times \mathrm{zscore}(f) \times \mathrm{sinal\_físico}(f),
\end{equation}
onde \(\mathrm{importância}(f)\) é carregada de \texttt{shap\_per\_class.json} (importância por classe predita, quando disponível) ou de \texttt{shap\_feature\_importance.json} (Gini global como \emph{fallback}); \(\mathrm{zscore}(f)\) mede a anormalidade do valor no ponto em relação à distribuição global do \emph{dataset}; e \(\mathrm{sinal\_físico}(f) \in \{-1, +1\}\) codifica o efeito esperado (por exemplo, \(+1\) para \texttt{Indice\_Seca}, \(-1\) para \texttt{Precipitacao\_ma90}). A predição é classificada como \emph{AUMENTA RISCO}, \emph{REDUZ RISCO} ou \emph{neutro}; o \emph{top-6 features} por \(|\mathrm{contrib}|\) é exibido no painel lateral.

\subsection{Interface}

O \emph{front-end} é organizado em \textbf{seis \emph{cards}} dentro de um painel lateral de \SI{420}{px}:

\begin{enumerate}
  \item \textbf{Previsão} --- \emph{badge} com a classe, barras de probabilidade por classe, métricas do modelo, fonte do clima, granularidade do \emph{lookup} e \emph{trace} da decisão (\emph{argmax} vs.\ \emph{thresholds} calibrados).
  \item \textbf{Por que esse risco?} --- \emph{top-6 features} com valor, faixa típica (\(p_{10}\)--\(p_{90}\) da distribuição global), \emph{z-score} e direção do impacto.
  \item \textbf{Dados climáticos atuais} --- precipitação, dias sem chuva, umidade NASA, temperatura climatológica, estação, período do dia.
  \item \textbf{Índices de seca e aridez} --- Índice de Seca, SPI 1m/3m/6m, KBDI \emph{proxy}, Aridez De Martonne, VPD \emph{proxy}, Dias\_Secos\_90d.
  \item \textbf{Precipitação acumulada e anomalia} --- acumulados 30/90/180/365~dias com narrativa textual comparando precipitação atual e a média histórica do município/mês.
  \item \textbf{Histórico de focos} --- focos NASA FIRMS em 7/30/90/180/365~dias, dias desde o último foco, FRP médio/máximo 7d, FRP no ponto agora.
\end{enumerate}

O mapa estático (versão original da Seção~4.6 do TCC~II) permanece disponível como artefato \emph{batch} (\texttt{scripts/gerar\_mapa.py}).

% ============================================================================
% (G) NOVA SECAO 4.7 -- PSEUDO-LABELING DE UMIDADE
% ============================================================================

\section{Imputação de \emph{Umidade} via \emph{pseudo-labeling}}
\label{sec:metodologia_pseudo_label}

Como discutido na Seção~\ref{subsec:metod_apis_externas}, a integração via NASA POWER produziu \num{170465} linhas com \texttt{Umidade} real (\(\approx \SI{18,3}{\percent}\) da base). Para mitigar essa cobertura parcial sem aguardar 13~dias adicionais de enriquecimento, foi implementada a estratégia de \textbf{\emph{pseudo-labeling}}~\cite{lee2013pseudo}, com precedente recente em imputação climática~\cite{lopezgarcia2024spatiotemporal}. O \emph{script} \texttt{scripts/pseudo\_label\_umidade.py} executa as etapas:

\begin{enumerate}
  \item Treina um \textbf{LightGBM regressor}~\cite{ke2017lightgbm} sobre os \(\approx \num{170000}\) registros com \texttt{Umidade} real, usando 15~\emph{features} (precipitação, dias sem chuva, lat/lon, sazonalidade cíclica \texttt{Mes\_sin/cos}, histórico de fogo) --- \emph{features} que existem para todas as linhas, com ou sem \texttt{Umidade} real.

  \item Avalia o erro real em \emph{holdout} 20\%: \textbf{RMSE = \SI{4,33}{\percent} RH, MAE = \SI{3,19}{\percent} RH, \(R^2 = 0{,}933\)}. O \(R^2\) alto é cientificamente justificável: a umidade tem alta autocorrelação espaço-sazonal com as \emph{features} usadas --- lat/lon definem o regime de monção, mês define a estação seca/chuvosa, precipitação recente define o estado higroscópico do ar.

  \item Refit no \SI{100}{\percent} dos \emph{labels} reais e imputa os \num{763489} registros sem \emph{label} (\emph{clip} [5, 100]\%).

  \item Salva o \emph{dataset} enriquecido em \texttt{base\_de\_dados\_umidade\_pseudo.csv} com coluna auxiliar \texttt{Umidade\_origem} \(\in\) \{\texttt{nasa\_power}, \texttt{pseudo\_label\_lgbm}\} para auditoria.
\end{enumerate}

\paragraph{Status no \emph{pipeline} final.} O \emph{dataset} \texttt{base\_de\_dados\_umidade\_pseudo.csv} está disponível mas \textbf{não é usado nos números finais reportados} no Capítulo~5. Razão: incluir \texttt{Umidade} (pseudo-rotulada ou real) exigiria retreino completo do \texttt{ensemble\_stacking\_gbm} (\(\approx \SI{85}{\minute}\) CPU) + revalidação das \emph{thresholds} + revalidação temporal --- fora do orçamento computacional desta entrega. O artefato fica registrado como \textbf{etapa preparatória} para a sequência natural do trabalho. Cautela registrada: \emph{pseudo-labels} devem ser interpretados considerando potencial viés de confirmação~\cite{arazo2020pseudo}.
